% Vietnamese translation for OI Public Library
% Source-File: IOI中国国家候选队论文/2003/王知昆/附件.doc
\documentclass[11pt]{article}
\usepackage{oipl-vn}

\title{Phụ lục: Nguyên đề của một số bài ví dụ}
\author{Wang Zhikun}
\date{2003}

\begin{document}
\maketitle

\origtitle{Appendix: original statements and programs for several examples}
\sourcepath{IOI China candidate team papers / 2003 / Wang Zhikun / appendix.doc}

\section{Nguyên đề của một số bài ví dụ}

\subsection{Bãi tắm cho bò}

\paragraph{Mô tả bài toán.}
Vì John xây hàng rào quanh nông trại, đàn bò nổi giận và sản lượng sữa giảm
mạnh. Để lấy lòng chúng, John quyết định xây một bãi tắm lớn trong nông trại.
Tuy nhiên đàn bò của John có một thói quen kỳ lạ: mỗi con bò phải cho sữa tại
một vị trí cố định trong nông trại, và dĩ nhiên bò không thể cho sữa trong bãi
tắm. Vì vậy John muốn bãi tắm được xây không phủ lên các điểm cho sữa đó. Lần
này anh lại phải nhờ Clevow giúp đỡ. Bạn có thể giúp Clevow không?

Nông trại của John và bãi tắm dự kiến đều là hình chữ nhật. Bãi tắm phải nằm
hoàn toàn trong nông trại, và các cạnh của bãi tắm phải song song hoặc trùng
với các cạnh của nông trại. Bãi tắm không được phủ bất kỳ điểm cho sữa nào,
nhưng một điểm cho sữa có thể nằm trên biên của bãi tắm.

Dĩ nhiên Clevow muốn diện tích bãi tắm lớn nhất có thể, nên nhiệm vụ của bạn
là tính diện tích lớn nhất của bãi tắm.

\paragraph{Tệp vào.}
Dòng đầu tiên của tệp vào chứa hai số nguyên \(L\) và \(W\), lần lượt là chiều
dài và chiều rộng của nông trại. Dòng thứ hai chứa một số nguyên \(n\), là số
lượng điểm cho sữa. \(n\) dòng tiếp theo, mỗi dòng chứa hai số nguyên \(x\) và
\(y\), biểu diễn tọa độ của một điểm cho sữa. Mọi điểm cho sữa đều nằm trong
nông trại, tức là \(0<x<L\), \(0<y<W\).

\paragraph{Tệp ra.}
Tệp ra chỉ gồm một dòng, chứa một số nguyên \(S\), là diện tích lớn nhất của
bãi tắm.

\paragraph{Ví dụ vào ra.}
\begin{center}
\begin{tabular}{p{0.42\linewidth}|p{0.42\linewidth}}
\texttt{happy.in} & \texttt{happy.out} \\
\hline
\begin{verbatim}
10 10
4
1 1
9 1
1 9
9 9
\end{verbatim}
&
\begin{verbatim}
80
\end{verbatim}
\end{tabular}
\end{center}

\paragraph{Ràng buộc.}
\[
  0<n<5000,\qquad 1<L,W<30000.
\]

\subsection{Candy Box}

\paragraph{Mô tả bài toán.}
Cho một hộp kẹo được chia thành \(n\times m\) ô. Ô ở hàng \(i\), cột \(j\) có
\(a[i][j]\) viên kẹo. Ban đầu tenshi định tặng hộp kẹo này cho một PPMM nào
đó, nhưng ngay vào đêm trước khi tặng, một con chuột cực kỳ đáng ghét đã đột
kích hộp kẹo: một số ô bị cướp sạch và bị đục thủng. tenshi phải nhanh chóng
cắt ra từ hộp kẹo này một hộp kẹo hình chữ nhật; hộp mới không được có lỗ, và
tenshi muốn tổng số kẹo giữ lại trong hộp mới là nhiều nhất có thể.

\paragraph{Nhiệm vụ.}
Hãy giúp tenshi thiết kế một chương trình tính số kẹo nhiều nhất có thể được
giữ lại trong hộp kẹo mới.

\paragraph{Định dạng vào.}
Dữ liệu được đọc từ tệp \texttt{CANDY.INP}. Dòng đầu tiên chứa hai số nguyên
\(n,m\). Số thứ \(j\) trên dòng \(i+1\) là \(a[i][j]\); nếu số này bằng \(0\)
thì ô tương ứng đã bị cướp và có lỗ. Trong đó:
\[
  1\le n,m\le 1000,\qquad 0\le a[i][j]\le 255.
\]

Chú ý: bài này cho \(16\) MB bộ nhớ và giới hạn thời gian là \(2\) giây.

\paragraph{Định dạng ra.}
Ghi số kẹo lớn nhất ra \texttt{CANDY.OUT}.

\paragraph{Ví dụ.}
\begin{center}
\begin{tabular}{p{0.42\linewidth}|p{0.42\linewidth}}
\texttt{CANDY.INP} & \texttt{CANDY.OUT} \\
\hline
\begin{verbatim}
3 4
1  2  3  4
5  0  6  3
10 3  4  0
\end{verbatim}
&
\begin{verbatim}
17
\end{verbatim}
\end{tabular}
\end{center}

\paragraph{Ghi chú.}
Hình chữ nhật sau có tổng số kẹo lớn nhất:
\begin{verbatim}
10  3  4
\end{verbatim}

\subsection{Big Barn}

Bài gốc là phát biểu chính thức tiếng Anh của bài \textit{Big Barn}.

\paragraph{A Special Treat.}
Farmer John wants to place a big square barn on his square farm. He hates to
cut down trees on his farm and wants to find a location for his barn that
enables him to build it only on land that is already clear of trees. For our
purposes, his land is divided into \(N \times N\) parcels. The input contains
a list of parcels that contain trees. Your job is to determine and report the
largest possible square barn that can be placed on his land without having to
clear away trees. The barn sides must be parallel to the horizontal or vertical
axis.

\paragraph{Example.}
Consider the following grid of Farmer John's land where \texttt{.} represents a
parcel with no trees and \texttt{\#} represents a parcel with trees:
\begin{verbatim}
          1 2 3 4 5 6 7 8
        1 . . . . . . . .
        2 . # . . . # . .
        3 . . . . . . . .
        4 . . . . . . . .
        5 . . . . . . . .
        6 . . # . . . . .
        7 . . . . . . . .
        8 . . . . . . . .
\end{verbatim}
The largest barn is \(5 \times 5\) and can be placed in either of two locations
in the lower right part of the grid.

\paragraph{Program name.}
\texttt{bigbrn}

\paragraph{Input format.}
Line 1: Two integers: \(N\) (\(1 \le N \le 1000\)), the number of parcels on a
side, and \(T\) (\(1 \le T \le 10{,}000\)), the number of parcels with trees.

Lines \(2\ldots T+1\): Two integers (\(1 \le\) each integer \(\le N\)), the row
and column of a tree parcel.

\paragraph{Sample input (file \texttt{bigbrn.in}).}
\begin{verbatim}
8 3
2 2
2 6
6 3
\end{verbatim}

\paragraph{Output format.}
The output file should consist of exactly one line, the maximum side length of
John's barn.

\paragraph{Sample output (file \texttt{bigbrn.out}).}
\begin{verbatim}
5
\end{verbatim}

\section{Chương trình cho một số bài ví dụ}

\subsection{Chương trình giải Bãi tắm cho bò bằng thuật toán thứ nhất}

\begin{verbatim}
program happy;
var
  f:text;
  x,y:array[1..5002] of longint;
  maxl,n,best,a,b,c,w,l,i,j,high,low:longint;

procedure sort(l,r:longint);
var
  i,j:longint;
begin
  i:=l+random(r-l+1);
  a:=x[i]; b:=y[i]; i:=l; j:=r;
  repeat
    while (x[i]<a) or ((x[i]=a) and (y[i]<b)) do i:=i+1;
    while (x[j]>a) or ((x[j]=a) and (y[j]>b)) do j:=j-1;
    if i<=j then
      begin
        c:=x[i]; x[i]:=x[j]; x[j]:=c;
        c:=y[i]; y[i]:=y[j]; y[j]:=c;
        inc(i); dec(j);
      end;
  until i>j;
  if j>l then sort(l,j);
  if i<r then sort(i,r);
end;

procedure sort_y(l,r:longint);
var
  i,j:longint;
begin
  a:=y[(l+r) div 2]; i:=l; j:=r;
  repeat
    while (y[i]<a) do i:=i+1;
    while (y[j]>a) do j:=j-1;
    if i<=j then
      begin
        c:=y[i]; y[i]:=y[j]; y[j]:=c;
        inc(i); dec(j);
      end;
  until i>j;
  if j>l then sort_y(l,j);
  if i<r then sort_y(i,r);
end;

procedure max(a:longint);
begin
  if a>best then best:=a;
end;

begin
  assign(f,'happy.in');
  reset(f);
  readln(f,l,w);
  readln(f,n);
  for i:=1 to n do
    readln(f,x[i],y[i]);
  close(f);
  inc(n); x[n]:=l; y[n]:=0;
  inc(n); x[n]:=0; y[n]:=w;
  sort(1,n);
  best:=0;
  for i:=1 to n do
    begin
      high:=w; low:=0; maxl:=l-x[i];
      for j:=i+1 to n do
        if (y[j]<=high) and (y[j]>=low) then
          begin
            if maxl*(high-low)<=best then break;
            max((x[j]-x[i])*(high-low));
            if y[j]=y[i] then break
            else if y[j]>y[i] then
              if y[j]<high then high:=y[j]
              else
            else if y[j]>low then low:=y[j];
          end;
      high:=w; low:=0; maxl:=l-x[i];
      for j:=i-1 downto 1 do
        if (y[j]<=high) and (y[j]>=low) then
          begin
            if maxl*(high-low)<=best then break;
            max((x[i]-x[j])*(high-low));
            if y[j]=y[i] then break
            else if y[j]>y[i] then
              if y[j]<high then high:=y[j]
              else
            else if y[j]>low then low:=y[j];
          end;
    end;
  sort_y(1,n);
  for i:=1 to n-1 do
    max((y[i+1]-y[i])*l);
  writeln(best);
end.
\end{verbatim}

\subsection{Chương trình của Candy}

\begin{verbatim}
program candy;
const
  maxn=1000;
var
  left,right,high:array[1..maxn] of longint;
  s:array[0..maxn,0..maxn] of longint;
  now,res,leftmost,rightmost,i,j,k,n,m:longint;
  f:text;
begin
  assign(f,'candy.in');
  reset(f);
  readln(f,n,m);
  fillchar(s,sizeof(s),0);
  for i:=1 to m do
    begin
      left[i]:=1; right[i]:=m; high[i]:=0;
    end;
  res:=0;
  for i:=1 to n do
    begin
      k:=0; leftmost:=1;
      for j:=1 to m do
        begin
          read(f,now); k:=k+now;
          s[i,j]:=s[i-1,j]+k;
          if now=0 then
            begin
              high[j]:=0; left[j]:=1; right[j]:=m;
              leftmost:=j+1;
            end
          else
            begin
              high[j]:=high[j]+1;
              if leftmost>left[j] then left[j]:=leftmost;
            end;
        end;
      rightmost:=m;
      for j:=m downto 1 do
        begin
          if high[j]=0 then
            begin
              rightmost:=j-1;
            end
          else
            begin
              if right[j]>rightmost then right[j]:=rightmost;
              now:=s[i,right[j]]+s[i-high[j],left[j]-1]-s[i-high[j],right[j]]-s[i,left[j]-1];
              if now>res then res:=now;
            end;
        end;
    end;
  writeln(res);
end.
\end{verbatim}

\subsection{Chương trình giải Big Barn bằng thuật toán thứ hai}

\begin{verbatim}
program BigBarn;
var
  d:array[1..1000,1..1000] of longint;
  height,left,right:array[1..1000] of longint;
  leftmost,rightmost,res,i,j,k,t,n:longint;
  f:text;
begin
  assign(f,'bigbrn.in');
  reset(f);
  readln(f,n,t);
  fillchar(d,sizeof(d),0);
  for i:=1 to t do
    begin
      readln(f,j,k);
      d[j,k]:=1;
    end;
  close(f);
  for i:=1 to n do
    begin
      height[i]:=0; left[i]:=1; right[i]:=n;
    end;
  res:=0;
  for i:=1 to n do
    begin
      leftmost:=1;
      for j:=1 to n do
        if d[i,j]=1 then
          begin
            height[j]:=0; left[j]:=1; right[j]:=n;
            leftmost:=j+1;
          end
        else
          begin
            height[j]:=height[j]+1;
            if leftmost>left[j] then left[j]:=leftmost;
          end;
      rightmost:=n;
      for j:=n downto 1 do
        if d[i,j]=1 then rightmost:=j-1
        else
          begin
            if rightmost<right[j] then right[j]:=rightmost;
            k:=height[j];
            if right[j]-left[j]+1<k then k:=right[j]-left[j]+1;
            if k>res then res:=k;
          end;
    end;
  assign(f,'bigbrn.out');
  rewrite(f);
  writeln(f,res);
  close(f);
end.
\end{verbatim}

\end{document}
